\documentclass[11pt]{article}
\usepackage[a4paper,margin=1.9cm]{geometry}
\usepackage[T1]{fontenc}
\usepackage{textcomp}
\usepackage[spanish]{babel}
\usepackage{amsmath,amssymb}
\usepackage{lmodern}
\renewcommand{\familydefault}{\sfdefault}
\usepackage{enumitem}

\newcommand{\vect}[2]{\begin{pmatrix}#1\\#2\end{pmatrix}}
\newcommand{\vectiii}[3]{\begin{pmatrix}#1\\#2\\#3\end{pmatrix}}
\usepackage{tikz}
\usetikzlibrary{calc,arrows.meta}
\usepackage[table]{xcolor}
\usepackage{parskip}

\definecolor{unalblue}{RGB}{31,73,124}
\newlist{opciones}{enumerate}{1}
\setlist[opciones]{label=\alph*),itemsep=4pt,topsep=4pt,leftmargin=1.8em,font=\normalfont}
\setlist[enumerate,1]{itemsep=1em,topsep=0.8em}

\newcommand{\examfooter}{\vfill\rule{\linewidth}{0.4pt}\\[1pt]{\scriptsize\textit{Transcripción digital --- MathUNAL}\hfill\textcolor{unalblue}{\textbf{\thepage}}}}
\pagestyle{empty}
\newcommand{\nota}[1]{{\small\itshape\color{unalblue!70!black} #1}}

\begin{document}
\noindent
\begin{minipage}[t]{0.6\linewidth}
{\small\bfseries Geometría Vectorial y Analítica --- Parcial 3}\\
{\small Semestre 2021--02}
\end{minipage}%
\hfill
\begin{minipage}[t]{0.35\linewidth}
\raggedleft\small Escuela de Matemáticas. Universidad Nacional de Colombia, Sede Medellín.
\end{minipage}\\[2pt]
\rule{\linewidth}{0.6pt}

\vspace{4pt}
\nota{Nota: el examen fuente se titula "Examen Final -- Febrero 2022" (cubre todo el semestre 2021-02), pero está ubicado en el catálogo del archivo original bajo Tercer Parcial 2021-2; se transcribe tal cual, respetando el título real.}

\begin{center}
\textsc{Geometría Vectorial y Analítica}\\
Semestre 2021-02 --- Examen Final --- Febrero 2022
\end{center}

\begin{center}
\begin{tabular}{|c|c|c|c|c|c|c|}
\hline
Pregunta & 1 & 2 & 3 & 4 & 5 & 6 \quad 7\\
\hline
Puntaje & 12\% & 9\% & 15\% & 12\% & 24\% & 12\% \quad 16\%\\
\hline
\end{tabular}
\end{center}

\textbf{Instrucciones:} La duración del examen es 2 horas y media. Verifique que sus datos de identificación son correctos. La prueba consta de siete preguntas. Debe consignar sus respuestas en la encuesta de Google Forms que le fue enviada por correo. Este examen es individual y no está permitido pedir ayuda o buscar las respuestas en internet.

\begin{enumerate}
\item Considere los puntos en la imagen y suponga que todas las líneas se cortan ortogonalmente. Seleccione todas las afirmaciones que son necesariamente verdaderas.

\begin{center}
\begin{tikzpicture}[scale=0.55]
\coordinate (O) at (0,0);
\coordinate (I) at (3.9,1.3);
\coordinate (E) at (-3.6,-1.2);
\coordinate (C) at (-0.6,1.8);
\coordinate (G) at (0.55,-1.65);
\coordinate (A) at (-3.9,0.7);
\coordinate (F) at (3.3,3.1);
\coordinate (Hp) at (4.6,-0.8);
\coordinate (B) at (-2.4,1.1);
\coordinate (D) at (0.9,-0.7);
\draw[thick] ($(A)!-0.1!(F)$) -- ($(F)!-0.1!(A)$);
\draw[thick] ($(E)!-0.1!(I)$) -- ($(I)!-0.1!(E)$);
\draw[thick] ($(G)!-0.1!(C)$) -- ($(C)!-0.15!(G)$);
\draw[thick] ($(Hp)!-0.1!(F)$) -- ($(F)!-0.15!(Hp)$);
\foreach \p/\lab/\pos in {O/O/below left,I/I/below right,E/E/below left,C/C/above,G/G/right,A/A/left,F/F/above,Hp/H/below right,B/B/above,D/D/right}{
  \fill (\p) circle (1.3pt); \node[\pos,font=\small] at (\p) {$\lab$};
}
\end{tikzpicture}
\end{center}

\begin{opciones}
\item Los vectores $A$ y $E$ son ortogonales.
\item Si el origen está en $A$, entonces la pareja $(E,B)$ está positivamente orientada.
\item Si el origen está en $I$, entonces $H\cdot G>0$.
\item Si el origen está en $E$, entonces $B$ y $H$ son linealmente dependientes.
\item Ninguna de las anteriores.
\end{opciones}

\item Sean $A,B,C,D$ vectores en $\mathbb{R}^2$ y $M$ el punto de intersección entre $\mathcal{L}_{AC}$ y $\mathcal{L}_{BD}$. Suponga que $B-A=\vect{6}{9}$, $C-A=\vect{8}{2}$, y $D-B=\vect{2}{-6}$. Encuentre $t$ tal que $M=A+t(C-A)$.

\item Suponga que $a,b,c,d\in\mathbb{R}$, con $a$ y $b$ diferentes de cero, y $P\in\mathbb{R}^2$. Considere la recta $\mathcal{L}$ con ecuación
$$x+by+c=0.$$
De las siguientes rectas, seleccione todas aquellas que necesariamente son perpendiculares a $\mathcal{L}$.

\begin{opciones}
\item $\left\{X\in\mathbb{R}^2:(X-P)\cdot\vect{b}{-1}=0\right\}$.
\item $\left\{X\in\mathbb{R}^2: X=P+t\vect{a/b}{a},\ t\in\mathbb{R}\right\}$.
\item $\left\{\vect{x}{y}\in\mathbb{R}^2: y-d=b(c-x)\right\}$.
\item $\left\{X\in\mathbb{R}^2: X=s\vect{1}{a+b}+t\vect{1}{a},\ s,t\in\mathbb{R},\ s+t=1\right\}$.
\item $\left\{\vect{x}{y}\in\mathbb{R}^2: -bx-y+d=0\right\}$.
\item Ninguna de las anteriores.
\end{opciones}

\item Determinar $r,s\in\mathbb{R}$ tales que $C=rA+sB$, donde $A,B,C$ son tres vectores en el plano que satisfacen las siguientes condiciones:
\begin{itemize}
\item el par $(A,B)$ está negativamente orientado, y el triángulo $\Delta(OAB)$ tiene área 5,
\item el par $(A,C)$ está positivamente orientado, y el triángulo $\Delta(OAC)$ tiene área 160,
\item el par $(B,C)$ está positivamente orientado, y el triángulo $\Delta(OBC)$ tiene área 100.
\end{itemize}

\item En la gráfica del presente problema se asume que los cuadrados pequeños miden $1\times 1$. En la figura (a) se grafican los vectores $A$ en color azul y $B$ en color rojo. En la figura (b) se presentan sus correspondientes imágenes a través de una transformación lineal $T:\mathbb{R}^2\to\mathbb{R}^2$, es decir $T(A)$ en azul y $T(B)$ en rojo. Los vértices del rombo $\mathcal{R}$ (café) de la figura (a) tienen coordenadas enteras.

\begin{center}
\begin{tikzpicture}[scale=0.42]
\draw[step=1,gray!30,thin] (-5,-5) grid (5,5);
\draw[-{Latex}] (0,-5.3) -- (0,5.3) node[above]{$y$};
\draw[-{Latex}] (-5.3,0) -- (5.3,0) node[right]{$x$};
\fill[orange!70] (-3,1) -- (-1,2) -- (1,1) -- (-1,0) -- cycle;
\draw[-{Latex},thick,blue] (0,0) -- (-2,-2) node[left]{$A$};
\draw[-{Latex},thick,red] (0,0) -- (1,-3) node[below]{$B$};
\end{tikzpicture}
\hspace{1cm}
\begin{tikzpicture}[scale=0.42]
\draw[step=1,gray!30,thin] (-5,-5) grid (5,5);
\draw[-{Latex}] (0,-5.3) -- (0,5.3) node[above]{$y$};
\draw[-{Latex}] (-5.3,0) -- (5.3,0) node[right]{$x$};
\draw[-{Latex},thick,blue] (0,0) -- (-2,-2) node[left]{$T(A)$};
\draw[-{Latex},thick,red] (0,0) -- (4,1) node[right]{$T(B)$};
\end{tikzpicture}
\end{center}

\begin{enumerate}[label=(\alph*)]
\item Suponga que la matriz de la transformación es $[T]=\begin{bmatrix}a&b\\c&d\end{bmatrix}$. Encuentre los valores de $a,b,c$ y $d$.
\item Los vértices del rombo $\mathcal{R}$ (café) de la figura (a), tienen coordenadas enteras. Encuentre el área de $T(\mathcal{R})$.
\item Considere las rectas
$$\mathcal{L}: -4x+7y+10=0,\qquad T(\mathcal{L}): y=mx+p,$$
donde $T(\mathcal{L})$ es la imagen de $\mathcal{L}$ bajo la transformación lineal $T$. Encuentre los coeficientes $m$ y $p$.
\end{enumerate}

\item Sea $(A,B,C)$ una terna ordenada de vectores en el espacio orientada positivamente. De las siguientes ternas, seleccione todas aquellas que son necesariamente positivamente orientadas.

\begin{opciones}
\item $(B\times C, 7C, -2B)$.
\item $((-2A+5C),B,C)$.
\item $(7C,6B,5A)$.
\item $(-2A,5B,-2C)$.
\item Ninguna de las anteriores.
\end{opciones}

\item Considere las rectas $\mathcal{L}_1$ y $\mathcal{L}_2$ cuyas ecuaciones paramétricas son:
$$\mathcal{L}_1:\begin{cases}x=2+(-3)t,\\ y=-1+(3)t,\\ z=4+(9)t,\end{cases} t\in\mathbb{R}\qquad \mathcal{L}_2:\begin{cases}x=2+s,\\ y=-1-3s,\\ z=4-s,\end{cases} s\in\mathbb{R}$$

Determine si cada enunciado es verdadero o falso:

\begin{enumerate}[label=(\alph*)]
\item La recta $\mathcal{L}_1$ intersecta el eje-$z$.
\item Las rectas $\mathcal{L}_1$ y $\mathcal{L}_2$ se intersectan en un único punto.
\item El vector $\vectiii{4}{1}{1}$ es normal a las rectas $\mathcal{L}_1$ y $\mathcal{L}_2$.
\item Hay un único plano que contiene las rectas $\mathcal{L}_1$ y $\mathcal{L}_2$.
\end{enumerate}

\end{enumerate}

\examfooter
\end{document}
